\documentclass[12pt]{article}
\usepackage[utf8]{inputenc}
\usepackage[T1]{fontenc}
\usepackage{hyperref}
\title{Yoba Language Resources for K-12 Curriculum Development}
\author{Compiled by Claude Sonnet 5.5 (claude-sonnet-5-5) for Andr\'as Kornai}
\date{}
\begin{document}
\maketitle

\section{Introduction}
Yoba (ISO 639-3: yob; Glottocode yoba1237) is an extinct Austronesian language of Central Province, Papua New Guinea. Glottolog places it in Oceanic, Papuan Tip, Central Papuan, Oumic, Magoric. Tom Dutton's notes place it north of Magori. The last speaker probably survived into the 1980s. Speakers shifted to neighbouring Papuan languages. The Wikipedia article relies largely on a single source (Dutton), so very little is documented.

\section{Primary Educational Resources}
None. No teaching materials of any kind were found.

\section{Linguistic and Reference Materials}
\begin{itemize}
\item Wikipedia, ``Yoba language''. \url{https://en.wikipedia.org/wiki/Yoba_language}
\item Max Planck EVA, numerals database entry for Yoba (Dutton's notes; the word for 5 means ``hand''). \url{https://lingweb.eva.mpg.de/channumerals/Yoba.htm}
\item Glottolog languoid page for Yoba. \url{https://glottolog.org/resource/languoid/id/yoba1237.bigmap.html}
\end{itemize}

\section{Religious and Community Resources}
None found.

\section{Digital and Audio Resources}
Wikipedia mentions that the Tom Dutton collection at Paradisec includes Yoba materials. A direct link was not found in the search results, so none is given. The collection should be checked directly.

\section{Curriculum Development Context}
Yoba is extinct and its documentation is minimal, so a language curriculum is not feasible from what the searches found. At most, a heritage module on the history and the Central Province language setting could be developed. The Paradisec materials should be investigated before anything else.

\end{document}